% ============================================================================
% 03-fundamentacao.tex — Fundamentação Teórica (módulo 4/8)
% ============================================================================
\chapter{Fundamentação Teórica}

\section{LLMs e raciocínio matemático}

Modelos de linguagem de grande escala, quando treinados em \textit{scale}
adequado, exibem capacidades preditivas descritas por leis de potência
\cite{hoffmann_chinchilla}. Modelos treinados em código demonstram
transferência de raciocínio estruturado \cite{chen_codex}, e técnicas de
elicitação de predições latentes estendem o aproveitamento dessas
capacidades \cite{ryoo_eliciting}. No raciocínio, \textit{Chain-of-Thought}
\cite{wei_cot} mostrou que instruir o modelo a produzir passos intermediários
melhora a exatidão aritmética em conjuntos como GSM8K \cite{cobbe_gsm8k}.
Agentes que interagem com ferramentas e refletem sobre falhas —
\textit{ReAct} \cite{yao_react} e \textit{Reflexion} \cite{shinn_reflexion} —
estendem o paradigma para tarefas interativas. O relatório técnico do GPT-4
\cite{openai_gpt4} documenta desempenho em exames padronizados.

Benchmarks matemáticos padronizados incluem MATH \cite{hendrycks_math}
(12.500 problemas de competição), GSM8K \cite{cobbe_gsm8k} (aritmética de
palavras) e MMLU \cite{hendrycks_mmlu} (conhecimento amplo). Para prova
formal assistida por máquina, LeanDojo \cite{yang_leandojo} integra LLMs ao
proverdor Lean. No domínio IMO, AlphaGeometry \cite{alpha_geometry}
combinou busca simbólica com heurística neural e provou 25/30 problemas de
geometria olimpíada; o marco de 2024 (AlphaProof) estendeu o resultado a
medalha de prata, ainda que sem publicação revisada por pares citável neste
relatório.

\section{Sistemas multiagentes}

A orquestração de múltiplos LLMs/agentes para tarefas complexas é tema
central de AutoGen \cite{wu_autogen} (conversas entre agentes programáveis) e
MetaGPT \cite{hong_metagpt} (papéis com SOPs). O OpenCode Ecosystem Core
implementa variante com contrato A2A via \textit{Blackboard}, memória
metacognitiva compartilhada (\textit{MetaBus}), \textit{Trust Engine} e
política SDD/TDD. A orquestração que conduziu este estudo usa o padrão
MASWOS (16 estágios) como \textit{scaffold} de raciocínio imposto ao LLM —
equivalente funcional a \textit{CoT}, mas com etapas explícitas de
verificação e conformidade.

\section{Crítica epistemológica e anti-overclaim}

Schaeffer e colaboradores \cite{schaeffer_mirage} argumentam que capacidades
aparentemente emergentes podem ser artefato de métricas descontínuas; a
seleção de métricas contínuas (acurácia por problema, \textit{score}
ponderado) e a replicação são, portanto, requisitos de rigor. A política do
repositório (CORRIGENDUM.md) registra alegações anteriores incorretamente
classificadas e exige validação externa para qualquer rótulo de excelência. O
presente relatório adota a mesma postura: nenhum resultado é declarado
\textit{superhuman}, \textit{A1} ou \textit{verificado} — apenas medido,
replicado quando possível e limitado explicitamente.

\section{Hipóteses de pesquisa}

\begin{enumerate}
  \item \textbf{H1}: modelos free reais, sob \textit{scaffold} MASWOS,
  resolvem problemas canônicos da IMO com acurácia superior à condição crua.
  \item \textbf{H2}: o \textit{score} composto (acurácia, velocidade,
  conformidade, disponibilidade) permite ordenar modelos free de forma
  reprodutível.
  \item \textbf{H3}: a orquestração multiagente (agentes de validação,
  solucionador determinístico anti-\textit{leak}, roteamento R500) é
  necessária para que o \textit{benchmark} seja auditável.
\end{enumerate}